\section{Introduction}\label{sec:introduction}
A data owner may need to demonstrate that public cluster centers give a reported assignment cost on private records. Examples include an analyst checking a deployed model on sensitive embeddings and an auditor requesting costs for several proposed model updates. Zero-knowledge proofs establish the specified computation without disclosing its private witness \cite{gmr1991knowledge,groth2016}. The difficulty is repeated high-dimensional arithmetic: a direct nearest-center audit evaluates distances in $d$ dimensions for every record and every candidate center. In our Groth16 implementation, a batch of eight 784-dimensional MNIST records requires 1.31 million constraints and 18.5 seconds of proving.

The same committed records can be queried many times. This motivates a two-stage construction. Stage~1 proves once that compact statistics were derived correctly from every committed original row. Stage~2 authenticates the requested statistics and proves nearest-center assignments and their aggregate cost. The central systems question is which statistics allow useful later queries while avoiding repeated original-space work.

For one fixed center set, the geometry is elementary. An orthonormal basis spanning the center differences preserves all assignment comparisons; adding each record's squared residual norm recovers every center distance. This follows from Pythagoras and belongs to the established projection-plus-offset approach to clustering sketches \cite{feldman2013tiny,feldman2020tiny,cohen2015dimensionality}. We use this identity as an interface to a proof system. We do not propose a new dimensionality-reduction identity or a new coreset theorem.

Fixed-center answers alone are insufficient to justify a reusable representation. A stronger task-specific baseline computes and proves the exact winning center and minimum cost in Stage~1, then authenticates those two values in later subset queries. It is cheaper for a single fixed center set. Exact affine scores provide another compact baseline with the same fixed-center information. Both measured implementations specialize preprocessing to the selected centers, so their Stage~1 work is repeated when those centers change.

We therefore study audits under changing public centers. We augment the deployed-center span with principal components of training residuals, and store the resulting head coordinates and residual energy. New centers enter Stage~2 through public head coordinates and residual offsets. They change the query inputs without changing the catalog or the query setup. This yields an exact real-arithmetic distance when the new centers remain in the represented affine span, and an approximation otherwise. The approximation drops a residual cross term, which we expose explicitly and evaluate against original-space distances.

The contribution has three parts:
\begin{enumerate}[leftmargin=*]
\item A checked integer, original-data-bound implementation of head and energy statistics, with all-row Stage~1 coverage and authenticated subset queries.
\item A reusable nearest-center circuit with public new-center statistics and a proved sum of compressed assignment margins. One setup serves center sets of the same shape; the margin is a diagnostic, not an original-space certificate.
\item Matched native and actual-proof measurements, including cached exact answers, affine scores, PCA with energy, direct Full proofs, and sparse-JL controls. Across three 64-/256-row proof workloads, reuse crosses the tested center-specific cached baseline at the second center set and reduces ten-set cumulative measured-component cost by factors of 6.82--7.76. Learned head and PCA with energy have similar proof cost; the learned head has higher assignment agreement on these proof catalogs. A separate eight-corpus, five-start campaign checks fixed-center fidelity.
\end{enumerate}

\section{Related work}\label{sec:related}
\paragraph{Projection and residual offsets}
Feldman, Schmidt and Sohler reduce clustering dimension while retaining an additive cost independent of the queried solution \cite{feldman2013tiny,feldman2020tiny}. Their arXiv extended version states the matrix Pythagorean identity (Claim~3) and, for centers in a subspace, the decomposition into projection cost and orthogonal residual cost (Eq.~9). Cohen et al. formalize projection-cost-preserving sketches with a solution-independent additive constant \cite{cohen2015dimensionality}. These results address approximation over classes of clustering solutions. Here we keep one energy value per authenticated record so that arbitrary requested subsets can be audited, and prove the deterministic derivation from the original commitment. The fixed-center identity is a special case of this classical geometry; the reuse circuit and its measured trade-offs are the focus.

\paragraph{Verifiable clustering and private analytics}
Prior work already combines private $k$-means and zero-knowledge techniques. Patel, Patel and Jinwala use secret sharing with a code-based identification scheme for horizontally partitioned clustering against malicious participants \cite{patel2013clustering}. Repetto et al. prove PAC-private mechanisms, including $k$-means, using RISC-Zero \cite{repetto2026pac}. Their objectives differ from auditing successive public center sets against an unchanged committed catalog. In nearby vector search, V3DB proves execution of published approximate-search semantics over committed snapshots \cite{v3db2026}. Our relation instead proves exhaustive nearest-center comparisons on authenticated statistics and reports the resulting approximation to original-space answers. These systems establish that neither ZK clustering nor audits over committed vector data are new categories.

\paragraph{Proof-oriented representations}
ezDPS studies a ZK inference pipeline including PCA, and Artemis develops commit-and-prove techniques for external data or model commitments \cite{ezdps2023,artemis2024}. Our PCA baseline proves an approved projected computation in the same backend; it does not reproduce ezDPS or prove PCA training. We use Groth16 over BN254 and Poseidon2 row-Merkle authentication \cite{groth2016,poseidon22023}. The contribution is a representation and circuit design within this setting, not a new proof primitive.

\paragraph{Privacy-preserving prototype models}
Schubert et al. implement a nearest-prototype learning model using fully homomorphic encryption and study the effects of integer encoding and distance computation \cite{schubert2026prototype}. Their objective is computation over encrypted data. Our objective is a zero-knowledge audit of specified nearest-center answers over committed records. Both settings make representation cost and integer semantics important, but encryption, proof soundness and approximation fidelity address different guarantees. The present experiments audit public centers; they do not prove private model training or neural feature extraction.

\paragraph{Data-oblivious alternatives}
JL and sparse JL preserve geometry for an independently fixed point set under their stated distribution and dimension assumptions \cite{johnson1984extensions,achlioptas2003database,kane2014sparser,makarychev2019performance}. A post-commit random map avoids requiring training data to choose a projection. We retain that protocol option and evaluate sparse-JL estimators, but the tested tails do not outperform PCA with energy on matched new-center cells. The binding guarantee for the learned representation comes from Stage~1 proofs, independently of whether a random tail is present.
